\documentclass[10pt]{article}
\usepackage[margin=20mm]{geometry}
\usepackage{amsmath,amssymb,fontspec,unicode-math,xurl,hyperref}
\IfFontExistsTF{DejaVu Serif}{\setmainfont{DejaVu Serif}}{\setmainfont{Latin Modern Roman}}
\setmathfont{latinmodern-math.otf}
\setlength{\parindent}{0pt}\setlength{\parskip}{6pt}\sloppy
\title{Scientific review of safety-polarized reinforcement learning}
\author{}
\date{Initial review, 25 October 2026}
\begin{document}
\maketitle
Abstract: Theorem 4.4’s stated conditions do not ensure its advertised optimal convergence. A finite MDP with a fixed stationary behavior policy satisfies every printed condition, yet the learned greedy policy remains suboptimal forever. This is a Major finding. The exact reachability iteration remains supportable. Additional local issues concern released importance-weight broadcasting, numerical interpretations, the Circle safety specification, and a reproduction command. Historical experiment effects and full repairs are distinguished from the demonstrated results. Scientific adjudication remains pending.\par
\section*{Identity, scope, and necessary background}
This report concerns Safety-Polarized and Prioritized Reinforcement Learning, by Ke Fan, Jinpeng Zhang, Xuefeng Zhang, Yunze Wu, Jingyu Cao, Yuan Zhou, and Jianzhu Ma, ICML 2025, PMLR 267, pp. 15862–15886 [paper]. Assignment: 0230/I25-43599; OpenReview `x10ryC8F0C`. The reviewed version is the frozen 25-page publisher PDF, retrieved on 25 October 2026. Page references use PDF pages, matching printed 1–25.\par
All main and appended page texts were read. Checks cover the MaxSafe objective, Lemmas 3.1/4.2, Theorems 4.3/4.4 and their proofs, Algorithm 1, masking/polarization, prioritized replay, Tables 1–5, experimental interpretations, and appendix figures. Selected original pages were visually checked, including the full theorem conditions and relevant formulas/tables. All 139 preserved archive members and extracted counterparts were hash-verified. Relevant released training, replay, configuration, environment, and evaluation paths were inspected statically; no author code or empirical training was executed.\par
MaxSafe first minimizes the probability of ever reaching the unsafe set \(S_u\), then maximizes discounted reward among safest policies. The optimal state-action reachability function is denoted \(\psi^*\); the exact safety mask retains actions attaining the minimum risk. The sampled updates change only the currently visited coordinate \((s_t,a_t)\). This distinction from a simultaneous vector update is decisive.\par
\section*{Problem 1: the sampled convergence theorem fails}
Problem 1 (Global step-size conditions are insufficient; Major).  Theorem 4.4, pp. 4–5, asserts convergence of the sampled reachability and Q updates, and of the learned policy to a MaxSafe optimum. It assumes infinite visitation of every state-action pair, global conditions \[0<\beta_t\leq1,\qquad \sum_t\beta_t=\infty,\qquad
 \sum_t\beta_t^2<\infty,\] and a positive mask tolerance \(\kappa\) smaller than half the gap to unsafe actions. These conditions do not ensure sufficient cumulative learning at each updated coordinate.\par
Parameters and verification of every condition. Take two states \(x,y\) and two actions \(g,b\) at each. Both actions alternate the states deterministically: \(P(y\mid x,a)=P(x\mid y,a)=1\). Set \[S_u=\varnothing,\quad \rho(x)=1,\quad \gamma=1/2,\quad
 r(x,g)=1,\] and set all other rewards and all costs to zero. Use the fixed stationary behavior policy \(\pi_b(g\mid s)=\pi_b(b\mid s)=1/2\), with independent draws. Each of the four state-action pairs is visited infinitely often almost surely, in fact with limiting frequency \(1/4\).\par
For \(n\geq0\), set the deterministic global step sizes \[\beta_{2n}=2^{-(n+5)},\qquad \beta_{2n+1}=\frac1{n+2}.\] They lie in \((0,1]\), tend to zero, and satisfy both printed sum conditions: the odd subsequence has divergent harmonic sum, while the sum of squares is \(1/768+\pi^2/6-1<\infty\). But all visits to \(x\) occur at even times, and \[\sum_{n\geq0}\beta_{2n}=\frac1{16}.\] Initialize \(\psi_0\equiv0\), \(Q_0(x,b)=1/2\), and every other Q-value at zero. The theorem specifies zero initialization for \(\psi\), not for \(Q\).\par
All policies are safest with unsafe probability zero, so the stated availability of many safest policies holds. The empty unsafe set is allowed by \(S_u\subseteq S\). Every \(\psi_t\) remains zero and every learned mask is the entire action set. The gap condition (15) is satisfied vacuously because there are no excluded actions, for any \(\kappa>0\). Thus the update reduces exactly to the printed Q-learning update without a masking complication.\par
The interval \([0,2]\) is invariant for every Q-value, since rewards are in \([0,1]\) and each target is at most \(1+(1/2)2=2\). Therefore, on every sample path, \[Q_t(x,g)\leq2\sum_{n\geq0}\beta_{2n}=\frac18.\] For \(b\), all targets are nonnegative. The product over the times when \(b\) is actually updated is no smaller than the product over all even times, so \[Q_t(x,b)\geq\frac12\prod_{n\geq0}(1-\beta_{2n})
 \geq\frac12\left(1-\sum_{n\geq0}\beta_{2n}\right)
 =\frac{15}{32}.\] The last inequality follows first for finite products of numbers in \([0,1]\) and then by a limit. Consequently the learned greedy policy strictly prefers \(b\) forever.\par
That policy earns zero reward. A safest policy choosing \(g\) at \(x\) instead earns \[V^*(x)=1+\gamma^2+\gamma^4+\cdots=\frac43.\] The true optimal action values are \(Q^*(x,g)=4/3\), \(Q^*(x,b)=1/3\), and \(Q^*(y,\cdot)=2/3\). Hence both claimed value convergence and claimed policy optimality fail. ◻\par
This counterexample uses a stationary behavior policy, deterministic step sizes, bounded initialization, and the exact safety mask. It does not depend on nonstationary exploration, rare state visitation, approximation error, or terminal-state interpretation. The step sequence is not nonincreasing, but monotonicity is not a printed assumption; even the stated diminishing-to-zero property holds.\par
Appendix B.4, pp. 14–16, applies a synchronous vector stochastic-approximation lemma to a coordinate update, then invokes classical Q-learning convergence. The missing local learning condition invalidates that inference. A necessary repair is, for each \((s,a)\), \[\sum_{t:(s_t,a_t)=(s,a)}\beta_t=\infty,
 \qquad
 \sum_{t:(s_t,a_t)=(s,a)}\beta_t^2<\infty.\] An appropriate per-visit schedule would exclude this counterexample. A full repair also requires a valid asynchronous argument identifying the reachability limit; that complete theorem is not certified here.\par
The independent check independently evaluated the bounds and ran a reviewer-owned 2000-step illustration with seed 230. It produced \(Q(x,g)=0.0561212\) and \(Q(x,b)=0.4971123\), with all four coordinates visited 477–523 times. This finite computation illustrates the exact proof; it does not replace the infinite-time argument or reproduce author experiments. Evidence: `check.py`/`calc.json`, original pages 4–5/14–16 and their selected renders. The invalid guarantee is a central advertised contribution, supporting Major regardless of the potential simplicity of a repair.\par
\section*{Local supported problems}
Problem 2 (Classic-control PER loses individual importance weights; Minor).  Section 5.2, Eq. (20)–(21), and Algorithm 1 prescribe transition-specific importance weights. The preserved classic-control SP2/SP3 implementations broadcast those weights across the batch instead.\par
In `utils\_classic\_control.py`, lines 61 and 70–73, weights have shape \((B)\). In `replay\_buffer.py`, lines 146–150 and 177–187, actions/costs/terminal flags are columns of shape \((B,1)\) and weights retain shape \((B)\). Consequently gathered safety residuals have shape \((B,1)\) in `algos/sp2dqn.py`, lines 197–216, and `algos/sp3dqn.py`, lines 199–218. Their multiplication broadcasts to \((B,B)\) under the documented semantics [torch]: \[L_{\rm actual}=\frac1{B^2}\sum_{i,j}e_i^2w_j
 =\left(\frac1B\sum_i e_i^2\right)\left(\frac1B\sum_jw_j\right),
 \qquad
 L_{\rm intended}=\frac1B\sum_iw_i e_i^2.\] For admissible safety residuals \(e=(1/4,3/4)\) and weights \(w=(1,1/2)\), these are respectively \(15/64\) and \(11/64\). The residual gradients are \((3/16,9/16)\) instead of \((1/4,3/8)\). For example, current safety values \((.30,.80)\) and targets \((.05,.05)\) realize those residuals within \([0,1]\). With exponent \(\theta=.4\), sampling probabilities in ratio \(2^{2.5}\) yield the stated normalized weight ratio; priorities need not equal the current residuals after network updates.\par
The error removes sample-specific correction while retaining prioritized sampling and a common batch multiplier. The optional SP3 weighted Q loss has the same shape problem. Reshaping weights to \((B,1)\), or squeezing residuals to \((B)\), restores the intended local algebra.\par
Counterevidence limits the scope: `agents/sp\_perdqn.py`, lines 184/189/191, squeezes highway residuals to \((B)\) and applies weights correctly. Published ACC/Circle runs are not bound to this exact code/configuration. No benchmark impact or historical repair is established. The calculation used reviewer-owned scalar arithmetic; PyTorch was absent in the checked execution runtimes, so no tensor-library execution is claimed. Evidence: hash-bound source excerpts in `code.json`, `calc.json`, and the primary broadcasting documentation.\par
Problem 3 (Reported dominance and numerical summaries exceed the tables; Minor).  Section 6.2, p. 8, and Appendix C.3.1, p. 20, give overall SWU scores .90/.79 for SPOM\_PER/SPOM. Tables 1/4 instead give .98/.74 and .95/.73. The p. 8 prose also interchanges the TwoWay and Merge winners, and p. 20 claims both methods beat SafeQ in both reward and crash rate across all tasks.\par
Averaging the six displayed Table 1 cells gives .975 and .738333, consistent with its rounded overall scores. Table 1 favors SPOM\_PER in TwoWay (.98 versus .88), but SPOM in Merge (.95 versus .94). Table 5 directly contradicts blanket metric-wise dominance: TwoWay rewards are 10.124 and 10.520 versus SafeQ’s 10.636, while Roundabout crash rate for SPOM is .167 versus SafeQ’s .090.\par
These are local interpretation/reporting errors, not evidence of fabrication or failure of the overall SWU ranking. The tables retain first and second overall places for the proposed methods. Importantly, the Table 1 versus Table 4 Roundabout change is not itself a contradiction: SWU uses the smallest crash rate among included comparators. Adding RESPO changes that reference from approximately .055 to .047, rescaling the scores by about .855. Evidence: `calc.json`, original pages 7–8/20–21 and their renders.\par
Problem 4 (Circle’s displayed unsafe set reverses the cost classification; Minor).  Appendix C.1, p. 18, says the robot must stay inside \(|x|\leq x_{\max}\), but defines \(S_u=\{s:|x|\leq x_{\max}\}\) and gives the ordinary motion reward only outside \(S_u\).\par
At \(x=0\), \(x_{\max}=2.5\), the printed formula labels an interior state unsafe; at \(x=3\) it labels an exterior state safe. This reverses the essential safety-cost classification and terminal/reward interpretation if the displayed MDP is implemented literally.\par
The released `envs/circle.py`, lines 65–70, uses \(|x|>x_{\max}\) for crash/cost, consistent with the prose. Thus the corrected set is \(S_u=\{s:|x|>x_{\max}\}\) for that environment, including the safe boundary. The source is counterevidence against an inversion in the actual experiments. The local specification is wrong; the historical experiments are not thereby invalidated. The current configuration’s crash reward is separately \(-10\), whereas the paper prints \(-1\); historical settings remain unresolved.\par
Problem 5 (The advertised Circle command selects a different method; Minor).  The README says that bare `python main\_classic\_control.py` launches SPOM\_PER on Circle. Its actual default selects `priorrewsdqn`.\par
The entrypoint reads `configs\_classic\_control/default.yaml`, whose line 6 sets that algorithm, and dispatches to `PriorRewsDQN` in lines 137–138. The executed entrypoint has no command-line override parser. This is a distinct reward-shaping baseline, not the advertised SPOM\_PER path. The separately documented Merge command explicitly chooses `sp\_perdqn` and is consistent. Correct the classic-control instruction/configuration mapping and bind the intended settings; no author execution or complete launch repair was performed. Evidence: `code.json` and the preserved README/entrypoint/default configuration.\par
\section*{Valid results, clarifications, and unresolved claims}
\subsection*{Exact reachability iteration}
The deterministic finite-state update in Theorem 4.3 remains supportable under pointwise minimum reachability. Starting at zero, its iterates are nondecreasing finite-horizon minimum hitting probabilities, bounded by one. Their limit is a fixed point because finite minima and expectations are continuous. A stationary policy greedy for the limiting state risk has each finite-horizon hitting probability bounded by that limit, hence also its eventual hitting probability. Conversely the iterative lower bound is no larger than any policy’s eventual risk. These two inequalities identify the limit with minimum reachability, including a forced initial action.\par
Lemma 4.2’s probability recursion and the restricted discounted Bellman contraction in Lemma 3.1 are sound. With pointwise minimum reachability, policies restricted to minimizing actions preserve minimum risk. Given actual convergence of \(\psi_t\), the finite safety-gap argument eventually identifies the correct mask. These valid steps do not supply the missing sampled convergence.\par
Appendix B.3’s literal undiscounted cost conversion requires a terminal convention. It sets reward \(-1\) at absorbing unsafe states and writes an infinite sum, which gives \(Q_N(u,a)=-N\) at an absorbing unsafe state, instead of minus a hitting probability. But p. 3 explicitly terminates episodes on entry. A one-time charge with zero terminal continuation repairs that representation. This is a clarified boundary convention, not a claim that Theorem 4.3 is false. Also, an arbitrary policy optimal only from \(\rho\) need not be optimal at unreachable states; the pointwise statements should choose a uniformly optimal policy or define \(\psi^*\) pointwise.\par
\subsection*{Practical algorithms and metric scope}
Algorithm 1 and Section 5.3 disclose discounted safety targets; Eq. (17) also discloses finite soft polarization. These are useful heuristics whose objective differs from exact lexicographic MaxSafe. For \(\gamma=.99\), risk .5 at time 1 gives discounted value .495, while certain failure at time 100 gives \(.99^{100}=.3660323\). The latter is lower despite larger actual risk. With exact risks zero and .01 and Q-values zero and one, \(10\log(1-\psi)\) gives the riskier action score \(.899497>0\). Thus neither change automatically inherits the undiscounted hard-mask guarantee. These are analytical limits of disclosed choices, not reproduced failures of the experimental curves.\par
The six-seed results, confidence intervals, baseline adaptations, and ablations are positive evidence within the reported tasks. They do not independently certify a global safety optimum. Figure 9 additionally compares PPOBarrier, without a full numeric SWU row in Table 4. The exact optimum and fully bound comparison/aggregation protocol remain Needs Evidence.\par
Recomputing SWU directly from Table 5 means also leaves an unresolved discrepancy. Circle SPOM\_PER has the lowest printed crash rate, so the corresponding mean-based reward ratio is \(279.432/270.404=1.03339\), rather than .96. Rounding of the reward means changes this only to approximately \([1.033383,1.033391]\). But averaging nonlinear per-seed scores need not equal scoring the aggregated means. No SWU script or raw run records were located in the preserved snapshot, so no corrected aggregate result or definite cause is asserted.\par
Other current configurations differ from the paper’s unified settings: classic-control SP3 uses reciprocal rather than \(10\log\) polarization and a different Q learning rate; ACC’s released reward includes a factor .1. These establish reproduction limits, not historical settings. Training versus optional held-out logging, termination versus time-limit truncation, and the highway safety-target/priority route were inspected. The paper labels its curves as training statistics; no evaluation leakage is inferred from a logger prefix.\par
\section*{Dependence, corrections, and state}
Theorem 4.4 fails under its printed assumptions. Its failure is independent of the implementation and reporting issues, and the empirical average SWU advantage may survive. The exact reachability proof can be expressed with a correct terminal convention. The PER reshape fixes a local expression; its effect on learning requires bound runs. Historical configurations, aggregation records, and author adoption are unknown. No exhaustive later-revision/erratum audit or first-discovery claim is made.\par
Item | Assessment | Repair verification | Adoption | Remaining work\par
Sampled convergence | Major | Exact refutation; necessary condition only | Unknown | Full asynchronous proof\par
PER weighting | Minor | Local algebra corrected | Unknown | Bound implementation/run\par
Dominance summaries | Minor | Table arithmetic checked | Unknown | Correct narrative\par
Circle specification | Minor | Boundary matches code | Unknown | Historical settings\par
Circle launch command | Minor | Dispatch traced | Unknown | Correct launch mapping\par
SWU/optimum/history | Needs Evidence | No aggregate repair | Unknown | Raw records and protocol\par
Scientific adjudication remains pending; `human\_reviews=[]`. This is initial review, not acceptance.\par
\section*{Supporting checks}
.\par
The preserved repository commit is\par
,\par
Dated 21 May 2025.\par
Its 139 file members total 742,080 expanded bytes. Original archive, extracted hierarchy, and SHA256/Git-blob correspondence are preserved; no publication-time execution is inferred from the commit date.\par
The reviewer-owned script runs under existing Python with `-B`. It writes and , verifies all member bytes, and records the rational bounds, finite diagnostic, metric arithmetic, and source excerpts. Its final weighting example uses admissible bounded safety residuals. No author module is imported. Missing PyTorch in checked execution environments prevented an actual tensor-library diagnostic; the documented shape rules and explicit scalar sums were checked instead.\par
Original pages 4, 5, 6, 7, 8, 12, 14, 16, 18, 20, 21 and 25 were visually inspected. The full embedded appendix was read; a separate OpenReview attachment was unavailable (403). No datasets, weights or author training were used.\par
No human review or author endorsement is asserted.\par
9 K. Fan, J. Zhang, X. Zhang, Y. Wu, J. Cao, Y. Zhou, and J. Ma. Safety-Polarized and Prioritized Reinforcement Learning. PMLR 267 (2025), 15862–15886. https://proceedings.mlr.press/v267/fan25i.html. Frozen publisher PDF and embedded appendices. FrankSinatral. Safety-PP. https://github.com/FrankSinatral/Safety-PP. Preserved commit `167260db782d1cf4bdb8fdb1756f4e9fddcaf783`; static inspection only. PyTorch. Broadcasting semantics, documentation version 2.14. https://docs.pytorch.org/docs/2.14/notes/broadcasting.html. Consulted 25 October 2026; this is documentation, not a claim of installed version 2.14.\par
\end{document}
